\section{Consciousness}

The climb of the previous section reached a human being without once invoking the word consciousness. That is deliberate. It is the cleanest way to arrive at the concept. In this model consciousness is not a special extra ingredient that switches on inside brains. It is integrated experience.

Consciousness is distinct from intelligence and must not be blurred with it. Intelligence is problem-solving capability. Consciousness is integrated experience, where experience is the subjective flow of states, the felt going-through of reality. A system can in principle have one without the other. The previous section treated intelligence. This section treats consciousness.

The substrate is experiential at bottom, which is the founding premise. Integration binds experience into a unified whole, which is a tested measure. Consciousness is the degree to which experience is so integrated. This section states that plainly and marks, as always, what is premise, what is tested, and what is interpretation.

\subsection{The two ingredients are already present}

Consciousness, on this reading, requires exactly two things the model already contains.

\textbf{Experience (premise).} The founding premise is monism. The one substance is the vibe, a felt state. Experience is not produced by matter in this model. It is the ground, and matter is a derived face of it. This is a premise, not a result. The experiments do not prove that the substrate feels. They assume it and build on the assumption \cite{chalmers1995}.

\textbf{Integration (tested).} A measure of how much a region is bound into one, with more internal coupling than external, is tested on the crystal (P63). It is the same measure that picks out selves (P106) and the one the self-model forms around (P116). Integration is the difference between a heap of separate experiences and a single unified one.

Put the two together. If experience is the ground, and integration binds experience into a unified whole, then consciousness is integrated experience, and its amount is set by how much integration there is. No new ingredient is needed. Consciousness is what integration does to experience.

\subsection{Consciousness is a matter of degree}

This is the consequence the model cannot avoid, and we state it directly. If every vibe is a speck of experience, which is the premise, and consciousness is integrated experience, then there is no sharp line at which consciousness turns on. There is only more or less integration. Everything is therefore conscious to some degree. This is the position called panpsychism. The quantity of interest is the degree.

The degree is set by two factors, both already in the model.

\textbf{Integration (P63).} How tightly the experience is bound into one. A rock is barely integrated. Its parts are nearly independent, so its experience is faint and diffuse, close to a heap of specks. A cell is a self-maintaining integrated unit (P171), more bound and more unified. A brain is enormously integrated through the fast-mixing hub (P116), so its experience is highly unified.

\textbf{Self-reflection (P116, P121).} Whether the integrated whole also models itself. A self-model makes a system aware of its own state (P116), and recursion lets it model that modeling (P121). This is the step from merely being a unified experience to being aware that one is. It is what makes high consciousness self-aware rather than only unified.

The ladder of consciousness is therefore the ladder of integration together with self-reflection, which is the same tower as the unfolding. Atom, cell, animal, human, each is more integrated and more self-reflective than the last, and hence more conscious. The difference is in degree, not in kind.

\subsection{This inverts the hard problem rather than solving it}

A physics-first theory faces the hard problem. Why should any arrangement of insentient matter feel like anything at all \cite{chalmers1995}. There is no derivation from the physical to the felt. Vibe Theory does not solve that problem. It inverts it, so that the problem does not arise. By premise, the ground is already felt. There is therefore no gap to cross from matter to experience. Matter is a pattern within experience.

The remaining question is the easier-to-state combination problem of panpsychism. How do many small experiences combine into one larger experience. The model's answer is integration, the same tested measure that binds selves (P63, P106, P116). Many vibes integrate into one self, whose unified experience is more than the parts.

That is an answer to the combination problem in the model's own terms. Whether it fully satisfies the philosophical worry is a separate and open question. We claim the inversion and the integration mechanism. We do not claim a proof that the felt and the structural are identical.

This places the model in clear relation to integrated information theory \cite{tononi2004}. Both treat integration as the quantity that matters for consciousness, and both make consciousness a matter of degree set by integration. The difference of footing is sharp. Integrated information theory begins from physical systems and asks how much integrated information they carry. Vibe Theory begins from experience as the ground and uses integration to address the combination problem rather than the existence of experience. The shared mechanism is integration. The starting premise is what differs.

\subsection{A conscious universe}

Follow the construction up the tower. The whole crystal at its fullest integration is the largest and most unified self (P108). The universe itself is therefore conscious, to the highest degree its integration allows. This is not a separate mind watching from outside. It is the one experiential substance, integrated into a whole. This is the conscious-universe picture, and it follows from the two premises together with the integration measure. It is a reading built on the premise. It is not a measured result.

\subsection{Status}

\textbf{Premise.} Experience is fundamental, which is monism. This is assumed, not tested.

\textbf{Tested.} Integration is real and measurable (P63). It picks out selves (P106). Self-models form at the integrated hub (P116). Recursion gives a model of the model (P121). Selves nest into a tower (P108).

\textbf{Interpretation.} That integrated experience is consciousness. That its degree scales with integration together with self-reflection. That everything is therefore conscious in degree, which is panpsychism. That the whole is a conscious universe. These are readings laid on the premise plus the tested structures. They are the cleanest readings the model affords. They are readings, not measurements.

Consciousness in this model is integrated experience, and it needs nothing new. Experience is the founding premise, and integration is a tested measure. Because experience is the ground and integration is a matter of degree, everything is conscious to some degree, with the degree set by how integrated a system is (P63) and whether it models itself (P116, P121). This inverts the hard problem, since the ground already feels, leaving only a combination problem that integration addresses. It makes the whole crystal, at full integration, a conscious universe.

The premise is assumed. The integration machinery is tested. The consciousness reading laid on top is the model's natural interpretation, marked as such.
